% !TeX root = euclid-without-words.tex
% !TeX Program=pdfLaTeX
%
% Intersecting tangent and secant
%
\begin{minipage}[t]{.48\textwidth}
\hspace{4em}
\begin{tikzpicture}[scale=.8]
% Circle goes from (a) to (b)
\coordinate (a) at (0,0);
\coordinate (b) at (6,0);
% Line containing lower points is (c) -- (d)
\coordinate (c) at (1,-3);
\coordinate (d) at (7,1.5);
% Point outside circle is (e)
\coordinate (e) at (1,5);
% Name the lower chord
\path [name path=chord] (c) -- (d);
% Draw a circle whose center is half-way between (a) and (b) through (a)
\coordinate (mid) at ($ (a)!.5!(b) $);
\node [draw,thick,circle through=(a),name path=circ] at (mid) {};
% Get intersection of the lower line with the circle
\path [name intersections={of=circ and chord,by={i1,i2}}];
% Draw the full secant
\draw [name path=secant2,thick,blue] (i2) -- node[right] {$\mathbf{a}$} (e);
% Get intersection of the secant with the circle
\path [name intersections={of=circ and secant2,by={s21,s22}}];
% Draw offset line from outside point with the intersection of the secant
\draw[thick,red]
  let \p1 = (s21), \p2 = (e) in
    (\x2-5pt,\y2) -- node[left] {$\mathbf{b}$} (\x1-5pt,\y1);
% Draw the tangent line
\coordinate (tan) at (tangent cs:node=circle, point={(e)}, solution=2);
\draw[thick,black!10!teal] (e) -- node[right,yshift=4pt] {$\mathbf{c}$} (tan);
% Display formula in color
\node at ($(mid)+(12pt,0)$) {$
\mathbf{\mathcolor{black!10!teal} c^2
= 
\mathcolor{blue} a \cdot
\mathcolor{red} b
}$};
\end{tikzpicture}
\end{minipage}
